\documentclass[10pt]{article}

\usepackage[margin=1in]{geometry}

\usepackage{amssymb}

\begin{document}

\section{Stochastic Context-Free Grammar (SCFG)}

\subsection{Notations}

Let $T=(V,E)$ be a rooted \emph{ordered} full binary tree.
$V_I$ is the set of internal nodes and $V_E=V\setminus V_I$ the set of leaves.
$E:V_I\to V\times V$ finds the two ordered children of an internal node.
We write $v\to(u,w)$ if $v$ is the parent of $u$ and $w$.
In this case, ${\rm par}(u)={\rm par}(w)=v$.
${\rm root}(T)$ is the root of the tree.
$T_v$ denotes the subtree descended from $v$.
For simplicity, $u\in T_v$ if $u\in V(T_v)$.

$\mathcal{G}(T,\Sigma)=(V_I\times\Sigma\cup\{S_0\},V_E\times\Sigma,\mathcal{R},\mathcal{P})$ is the SCFG derived from $T$
where $\Sigma$ is the alphabet of residues (nucleotides or amino acids), $S_0$ is the start state,
$V_I\times\Sigma$ consists of the non-terminal states
and $V_E\times\Sigma$ consists of terminal states.
$$
\mathcal{R}=\{[v,a]\Rightarrow([u,b],[w,c])|\mbox{$v\to(u,w)\in E(T)$ and $a,b,c\in\Sigma$}\}
$$
is the set of production rules (or simply rules) and $\mathcal{P}:\mathcal{R}\to\Re$ sets the rule probabilities.
If $[v,a]\Rightarrow([u,b],[w,c])$, its probability takes the form of $p_u(b|a,\theta)p_w(c|a,\theta)$,
parameterized by $\theta$.
$p_u(b|a)$ is the transition probability from $v={\rm par}(u)$ to $u$.
At the root, the probability of $S_0\Rightarrow[{\rm root}(T),a]$ is $q(a)$.

Intuitively, the SCFG generates residues at leaves.
Each derivation corresponds to one way to generate the residues.
Note that although the deduction below is based on the SCFG theory, the SCFG itself is implicit.
We do not need to explicitly construct the SCFG during implementation.

\subsection{The inside algorithm}

In the following, we assume $v\to(u,w)$.
$\sigma(v)$ denotes resides at leaves descended from $v$.

\begin{eqnarray*}
\alpha(v,a)&\triangleq&\Pr\{[v,a]\stackrel{*}{\Rightarrow}\sigma(v)\}\\
&=&\sum_{b,c}\Pr\big\{[v,a]\Rightarrow([u,b],[w,c])\stackrel{*}{\Rightarrow}\sigma(u)\sigma(w)\big\}\\
&=&\sum_{b,c}p_u(b|a)p_w(c|a)\alpha(u,b)\alpha(w,c)\\
&=&\left(\sum_b p_u(b|a)\alpha(u,b)\right)\cdot\left(\sum_c p_w(c|a)\alpha(w,c)\right)\\
&=&\alpha'(u,a)\alpha'(w,a)
\end{eqnarray*}
where
$$ \alpha'(u,a)\triangleq\sum_b p_u(b|a)\alpha(u,b) $$
$\alpha(S_0)=\Pr\{S_0\stackrel{*}{\Rightarrow}\sigma({\rm root}(T))\}$ is the probability of the tree.
The inside algorithm is equivalent to Felsenstein's algorithm.
The time complexity is $O(|V|\cdot|\Sigma|^2)$.

\subsection{The outside algorithm}

\begin{eqnarray*}
\beta(u,b)
&\triangleq&\Pr\{S_0\stackrel{*}{\Rightarrow}\sigma_1[u,b]\sigma_2\}\\
&=&\sum_{a,c}\Pr\big\{S_0\stackrel{*}{\Rightarrow}\sigma'_1[v,a]\sigma'_2\Rightarrow\sigma'_1([u,b],[w,c])\sigma'_2\stackrel{*}{\Rightarrow}\sigma'_1([u,b],\sigma(w))\sigma'_2\big\}\\
&=&\sum_{a,c}\beta(v,a)p_u(b|a)p_w(c|a)\alpha(w,c)\\
&=&\sum_ap_u(b|a)\beta(v,a)\left(\sum_c p_w(c|a)\alpha(w,c)\right)\\
&=&\sum_ap_u(b|a)\beta(v,a)\alpha'(w,a)
\end{eqnarray*}
For $v={\rm root}(T)$,
$$ \beta(v,a)=\Pr\{S_0\Rightarrow[v,a]\}=q(a) $$
where $q(a)$ is the ancestral state of the root. The probability of the tree
can be computed with
\begin{eqnarray*}
P(T)&=&\Pr\{S_0\stackrel{*}{\Rightarrow}\sigma({\rm root}(T))\}\\
&=&\sum_b\Pr\{S_0\stackrel{*}{\Rightarrow}\sigma_1[u,b]\sigma_2\stackrel{*}{\Rightarrow}\sigma_1\sigma(u)\sigma_2\}\\
&=&\sum_b\alpha(u,b)\beta(u,b)
\end{eqnarray*}

\subsection{The branch form}

We call $\alpha(v,a)$ and $\beta(v,a)$ the node forms.
We introduce the branch form in preparation for the next section.
It is also essential for fast nearest neighbor interchange (NNI) in rooted trees.

\begin{eqnarray*}
P(T)
&=&\sum_b\alpha(u,b)\beta(u,b)\\
&=&\sum_{a,b,c}\alpha(u,b)\beta(v,a)p_u(b|a)p_w(c|a)\alpha(w,c)\\
&=&\sum_{a,b}\alpha(u,b)p_u(b|a)\beta(v,a)\alpha'(w,a)
\end{eqnarray*}
If $u\not={\rm root}(T)$, let
$$
\eta_u(b|a)\triangleq\alpha(u,b)\beta(v,a)\alpha'(w,a)
$$
Then
$$ P(T)=\sum_{a,b}p_u(b|a)\eta_u(b|a) $$
When $u$ is an internal node, suppose $u\to(x,y)$. We have
$$
\eta_u(b|a)=\beta(v,a)\cdot\alpha'(x,b)\alpha'(y,b)\alpha'(w,a)
$$
With this equation, if we swap the sibling of $u$ with a child of $u$,
$\beta(v,a)$ and all three $\alpha'()$ stay the same.
We can update $\eta_u(b|a)$ without descending into the subtrees of $x$, $y$ and $w$,
and thus we can recompute the probability of the new tree in $O(|\Sigma|^2)$ time.
This gives us an efficient way to explore alternative tree topologies locally,
which is effectively NNI.

Note that we can also extend the definition $\eta_u(b|a)$ to the root by defining
$p_u(b|\cdot)=q(b)$ and $\eta_u(b|\cdot)=\alpha(u,b)$.
We will use this convention in the following.

\subsection{The posteriors}

Let $X_{v,a}(\pi)$ be the number of times we see $a$ at node $v$ in a
derivation. On a single tree, $X_{v,a}$ is either 0 or 1. We can compute
\begin{eqnarray*}
x_{v,a}&=&\sum_{\pi}X_{v,a}(\pi)\Pr(\pi|T)\\
&=&\Pr\{S_0\Rightarrow\sigma_1[v,a]\sigma_2\Rightarrow\sigma_1\sigma(T_v)\sigma_2|S_0\Rightarrow\sigma(T)\}\\
&=&\frac{1}{P(T)}\alpha(v,a)\beta(v,a)\\
&=&\frac{\alpha(v,a)\beta(v,a)}{\sum_b\alpha(v,b)\beta(v,b)}
\end{eqnarray*}
Obviously, $\sum_a x_{v,a}=1$, which is expected. Similarly, let
$X_{u,ab}(\pi)$ be the number of times $a\to b$ changes are observed between
$u$ and its parent $v$.
\begin{eqnarray*}
x_{u,ab}
&=&\sum_{\pi}X_{u,ab}(\pi)\Pr(\pi|T)\\
&=&\sum_c\Pr\{S_0\Rightarrow\sigma_1[v,a]\sigma_2\Rightarrow\sigma'_1([u,b],[w,c])\sigma'_2\Rightarrow\sigma'_1(\sigma(T_u),\sigma(T_w))\sigma'_2|S_0\Rightarrow\sigma(T)\}\\
&=&\frac{1}{P(T)}\sum_c\beta(v,a)p_u(b|a)p_w(c|a)\alpha(u,b)\alpha(w,c)\\
&=&\frac{\beta(v,a)p_u(b|a)\alpha(u,b)\alpha'(w,a)}{\sum_{a',b'}p_u(b'|a')\eta_u(b'|a')}\\
&=&\frac{p_u(b|a)\eta_u(b|a)}{\sum_{a',b'}p_u(b'|a')\eta_u(b'|a')}
\end{eqnarray*}
with $\sum_{a,b}x_{u,ab}=1$.

\subsection{Parameter estimation by EM}

The EM $Q$ function is:
$$
Q(\theta|\theta^t)=\sum_i\sum_u\sum_{a,b}x^{(i)}_{u,ab}\log p_u(b|a)+\sum_i\sum_c x^{(i)}_{{\rm root}(T),c}\log q(c)
$$
where $x^{(i)}_{u,ab}$ corresponds to the posterio at column $i$ in a multi-sequence alignment (MSA).
If we parameterize the SCFG by $p_u(b|a)$ with all transitions being free parameters,
the EM estimate of $p_u(b|a)$ is
$$
\hat{p}^{(t+1)}_u(b|a) = \frac{\sum_i x^{(i)}_{u,ab}}{\sum_i x^{(i)}_{{\rm par}(u),a}}
$$

\subsection{Numerical stability}

For numerical stability, we need to scale $\alpha$ and $\beta$.
We introduce $h(v)$ for $v\in T$ to scale $\alpha$.
To find a proper form of $h(v)$, define
$$ \tilde\alpha(v,a) \triangleq\frac{\alpha(v,a)}{\prod_{v'\in T_v}h(v')} $$
$$ \tilde\beta(u,b)  \triangleq\frac{\beta(u,b)}{h(u)\prod_{u'\in T\setminus T_u}h(u')} $$
$$ \tilde\alpha'(v,a)\triangleq\frac{\alpha'(v,a)}{\prod_{v'\in T_v}h(v')}=\sum_b p_u(b|a)\tilde\alpha(u,b) $$
$$ \tilde\eta_u(b|a) \triangleq\frac{\eta_u(b|a)}{\prod_{u'\not=u}h(u')}
= \tilde\beta(v,a)\cdot\tilde\alpha'(x,b)\tilde\alpha'(y,b)\tilde\alpha'(w,a)$$
We can prove
$$ \tilde\alpha(v,a)=\frac{1}{h(v)}\sum_{b,c}p_u(b|a)p_w(c|a)\tilde\alpha(u,b)\tilde\alpha(w,c)=\frac{1}{h(v)}\tilde\alpha'(u,b)\tilde\alpha'(w,c) $$
$$ \tilde\beta(u,b)=\frac{1}{h(u)}\sum_{a,c}p_u(b|a)p_w(c|a)\tilde\beta(v,a)\tilde\alpha(w,c)=\frac{1}{h(u)}\sum_ap_u(b|a)\tilde\beta(v,a)\tilde\alpha'(w,a) $$
\begin{eqnarray*}
P(T)&=&h(v)\prod_{v'}h(v')\cdot\sum_a \tilde\alpha(v,a)\tilde\beta(v,a) \\
&=&\prod_{v'\not=u}h(v')\sum_{a,b}p_u(b|a)\tilde\beta(v,a)\cdot\tilde\alpha'(x,b)\tilde\alpha'(y,b)\tilde\alpha'(w,a)
\end{eqnarray*}
There are more than one ways to define $h(v)$. We may choose:
$$ h(v)=\sum_{a,b,c}p_u(b|a)p_w(c|a)\tilde\alpha(u,b)\tilde\alpha(w,c) $$
This way, $\sum_a\tilde\alpha(v,a)=1$.

%\section{Emission Probability}
%
%$$\left(\begin{array}{lll}
%p & q & 0 \\
%p & q/2 & q/2 \\
%\end{array}\right)\Rightarrow
%\left(\begin{array}{lll}
%1 & 1 & 0 \\
%1 & 1/2 & 1/2 \\
%\end{array}\right)\Rightarrow
%\left(\begin{array}{lll}
%1 & 1-\epsilon/2 & \epsilon/2 \\
%1 & 1/2 & 1/2 \\
%\end{array}\right)
%$$
%$$\left(\begin{array}{llll}
%p^2 & 1-p^2 & 0 & 0 \\
%p^2 & pq & q^2 & pq \\
%\end{array}\right)\Rightarrow
%\left(\begin{array}{llll}
%1 & 1 & 0 & 0 \\
%1 & \frac{p}{1+p} & \frac{q}{1+p} & \frac{p}{1+p} \\
%\end{array}\right)\Rightarrow
%\left(\begin{array}{llll}
%1 & 1-\frac{\epsilon}{1+p} & \frac{\epsilon q}{1+p} & \frac{\epsilon p}{1+p} \\
%1 & \frac{p}{1+p} & \frac{q}{1+p} & \frac{p}{1+p} \\
%\end{array}\right)
%$$

\end{document}
